% !TeX root = ../papers/nc-p3.tex
\section{中心の次数安定化群と指数四の二重周期}
\label{sec:central-character-periods}
\begin{theorem}\label{thm:central-character-periods}
中心有限な標準正則有限四辺で、頂点指数を $s$ とする。中心の次数は経路長を4で割る。
非頂点の閉中心指標における次数作用の安定化群を $\mu_h$ とすれば、$h>1$ なら $2h\le s$ である。
一般指標の安定化群は自明である。特に $s=4$ では射影中心上の $\mathcal O(2)$ は豊富な直線束である。
この場合の非自明安定化群は正確に次の指標に出て、いずれでも経路長8の中心系は基点を持たない。
\begin{enumerate}
\item 定理\ref{thm:index-four-quartic-basepoints}の逆符号の二次元対または薄い二対。
\item 各頂点四次元の単純一つで、それが向き付き二次曲面の最小周期8の列の押下げである場合。
\end{enumerate}
従って経路長8でも消える基点があれば、そこは滑らかなAzumaya指標で、四次元単純の次数安定化群は自明である。
\end{theorem}
\begin{proof}
各頂点次元 $d$ の単純 $S$ の一周倍率安定化群を $\mu_m$ とする。倍率は戻り辺だけに掛ければよい。
任意の非頂点閉指標ではある正次数中心値が非零なのでこれらの群は有限である。
原始 $m$ 乗根を実現する頂点基底 $g$ は $m$ 乗がスカラーであり、全体の倍率で $g^m=1$ にできる。
各頂点の固有空間を取ると、通常の辺は固有指標を保ち、戻り辺だけが一つ進める。
一つ非零なら次も非零である。矢の共通核は単純性で零なので、巡回して全ての指標が全頂点に出る。
薄い単純なら $m=1$ である。非薄単純なら全固有空間の次元は2以上である。一次元固有空間の前位置の
対応固有空間に非零ベクトルを置くとその全矢像が直線となり、有限グラフ上の直線の伝播に反する。
従って非薄なら $d\ge2m$。$\mu_h$ は一つの中心繊維の単純型を置換し、その中の $S$ の軌道長は $h/m$。
$m$ は確かに $h$ を割る。表現が同型なら中心指標も同じだからである。
非薄型を含めば、次数 $s$ の被約traceに出る頂点次元の総和から $s\ge d h/m\ge2h$。
薄い型だけで $h>1$ なら同一台点での倍率の軌道長は $h$ であり、その一台点だけでは
異倍率間の拡張消滅で中心繊維が分解する。従って二台点以上でやはり $s\ge2h$。

一般指標が自明安定化群を持つ理由も記す。任意次数の四つの斉次中心Noetherパラメータの
経路長を $d_i$ とすると、全環はその多項式環 $P$ 上自由で、相対双対は
$\operatorname{Hom}_P(B,P)=B(\sum d_i-4)$ である。定理\ref{thm:central-cut-center-dichotomy}の
両側双対計算で、四元の差のKoszul行列式の次数を $\sum d_i$ と読めばそのまま従う。
同所の巡回汎関数と被約traceの制限により、全中心の標準加群は圧縮次数で $\omega_Z=Z(-1)$ となる。
一般安定化群の位数 $g$ は $Z$ の非零斉次片の次数の最大公約数である。もし $g>1$ なら $Z$ も
$P$ も圧縮次数 $g$ の倍数だけにあり、有限自由パラメータ上のHomで計算する $\omega_Z$ も同じ剰余類だけにある。
$Z(-1)$ と両立しない。

$s=4$ なら $h=1,2$ だけである。各指標の非零値を与える斉次中心元の次数の最大公約数が $h$ だから、
有限個のそれらを逆にして次数 $h$ の可逆な斉次分数を作れる。よって次数2移動は局所的に自明化され、
$\mathcal O(2)$ は直線束、高い冪の豊富性から豊富である。$h=2$ のとき上の不等式の等号を読む。
薄いものなら逆符号の二対、非薄で $m=1$ なら $d=2$ の逆符号対のみである。$m=2$ なら
$d=4$ の一型で、全固有空間が二次元となる。その固有空間の列は全矢の像条件を持つ八周期表現で、
元の単純はその押下げである。向き付き曲面列が四周期なら、各二次元辺の基底と接続直線の一意性から
四周期行列の二回反復に各辺のスカラーを掛けた形になる。八個のスカラーの積の平方根を四辺の一周倍率に選べば
頂点基底でそれを揃えられ、押下げは二次元単純二つに分解する。実際戻りの二層交換のスカラー行列を
対角化すればよい。よって最小周期は8。逆は系\ref{cor:quadrics-period-fold-simple}による。

最後の場合の八文字traceは二つの固有層を始点とする八周期二次元traceの和である。
どちらのテンソルも、その射影類から各位置の向き付き二次曲面を復元できる。
定理\ref{thm:index-four-quartic-basepoints}の隣接文字の収縮階数による証明で、残りの六辺の積空間も
全二行二列環であることを使えばよい。従って二つが比例すれば四段移動で曲面列が等しくなり、最小周期8に反する。
ゆえに和は非零である。残る二型の八次非消滅は同じ定理で証明した。なお八次が全て消えれば
四次中心元の平方も消えるため四次基点である。これで最後の主張も従う。
\end{proof}

\begin{corollary}\label{cor:index-four-no-divisorial-inertia}
十五次元成分閉包の有限四辺の正則整域窓では、射影中心 $X$ で位数二の安定化群が現れる軌跡の
余次元は2以上である。$H$ を次数1に対応する反射的因子類とすると $2H$ は豊富なCartier因子で
$K_X=-H$ である。
\end{corollary}
\begin{proof}
この閉包の中心Hilbert関数は定理\ref{thm:finite-domain-specialization-parameters}により一般四次超曲面のものに
等しく、十分大きい全整数で同じ三次多項式 $P(n)$ になる。$A=\mathcal O_X(2)$ とし、
$A^b$ の二つの一般切断で滑らかな曲線 $C$ を取る。中心は正規で、余次元2以上の集合は避けられる。
$L=\widetilde{Z(1)}$ は一般安定化群が自明なので階数1で、曲線に沿って直線束である。
実際次数2のVeronese上の奇成分は捩れがなく、余次元1の離散付値環上で自由となる。
乗法 $L^{\otimes2}|_C\to A|_C$ は生成点で非零である。指標で次数1の単元があることは
奇次数の斉次分数の単元があることと同値なので、一般の位数二の因子ではこの積が零点を持つ。
適切な偶次数の単元でずらしても、非零な奇次数値自体がそこにないからである。

十分大きい $n$ で $L\otimes A^n,A^n$ の大域切断の次元は $P(2n+1),P(2n)$ である。
二切断で制限すると曲線上のEuler数の差は、それらの差に幅 $b$ の後退差分を二度施したものになる。
従ってRiemann--Rochにより $\deg L|_C$ は $\deg A|_C$ の半分である。
実際 $P(2n+1)-P(2n)$ の二次最高係数は $P(2n+2)-P(2n)$ の半分である。
よって上の積の零点の次数は零となり、位数二の因子はない。余次元1で次数1単元が使えるため、
錐の非頂点部からの $\mathbb G_m$ 商はそこで通常の主束である。標準加群 $Z(-1)$ の層化から
$K_X=-H$ が余次元1で従い、両辺の反射的延長で全体の因子類として成立する。
\end{proof}
